We propose a doxastic reading of Propositional Lax Logic (PLL), the
intuitionistic modal logic of Fairtlough and Mendler whose single
modality \ensuremath{\bigcirc} has a flavour both of possibility and of necessity: we read
\ensuremath{\bigcirc}M as "M is believed", by a single idealised believer whose belief is
disciplined by evidence. The reading is defended by a chain of
machine-checked results. On the negative side, over classical logic the
reading is empty: on a Boolean algebra every nucleus is closed, so a
classical belief modality collapses to \ensuremath{\bigcirc}M \ensuremath{\equiv} M \ensuremath{\vee} a for a fixed
proposition a. Over intuitionistic logic belief is rich, and its
richness is evidential in a precise sense inherited from the constraint
reading.

The paper\textquotesingle s technical core is a new realisability semantics for PLL
over the constraint models of the original paper, in which atoms carry
evidence from a partial applicative structure and the clauses for \ensuremath{\supset} and
\ensuremath{\bigcirc} differ in how much a realiser is told about the future of the
inquiry. Three clause families arise --- uniform, strategy, and
presented-strategy --- and a sequence of separation theorems shows the
choice determines the logic: rigid evidence validates schemes PLL does
not prove; strategy evidence restores the countermodels but provably
cannot support a completeness theorem, by an obstruction created by the
modality\textquotesingle s own strategies; presenting the future at implications yields
completeness. The countermodels are supplied by a finitisation of the
original canonical-model construction in which Zorn\textquotesingle s lemma is replaced
by an iterated appeal to PLL\textquotesingle s own mechanised decision procedure, and
the resulting models are executable: a verified checker certifies them
by computation, and each world reads as a belief state --- what is
believed, and what is promised never to be evidenced.

Every result is formalised in Lean 4; the axiom footprint of each
theorem is tracked, and the greater part of the development, including
the entire decoration argument, is choice-free. We conclude with an
account of the method, in which several of the paper\textquotesingle s results were
found by the failure, itself provable, of the theorems we first
attempted.
